\documentclass[10pt,letterpaper]{article}
\usepackage[T1]{fontenc}
\usepackage{lmodern,amsmath,amssymb,booktabs,array,geometry,xcolor}
\usepackage[hidelinks]{hyperref}
\geometry{margin=0.76in,headheight=14pt,headsep=15pt}
\definecolor{ink}{HTML}{173A51}
\makeatletter
\def\ps@draft{\def\@oddhead{\small\color{ink}Spin(10) and global PQ: discussion draft\hfil 17 October 2026}
\def\@evenhead{\@oddhead}
\def\@oddfoot{\small Discussion draft; bounded results and open physical questions\hfil\thepage}
\def\@evenfoot{\@oddfoot}}
\makeatother
\pagestyle{draft}
\setlength{\parindent}{0pt}\setlength{\parskip}{6pt}
\newcommand{\PQ}{\mathrm{PQ}}
\newcommand{\Spin}{\operatorname{Spin}}
\newcommand{\Tr}{\operatorname{Tr}}
\newcommand{\hc}{\mathrm{h.c.}}
\newcommand{\pathref}[1]{\texttt{\detokenize{#1}}}
\newcommand{\sect}[1]{\vspace{5pt}{\large\bfseries\color{ink}#1}\par}
\begin{document}
\thispagestyle{draft}
{\LARGE\bfseries\color{ink} An AI-assisted investigation of a\par
Spin(10) $\times$ global U(1)$_{\PQ}$\par
gauge/scalar candidate\par}
\vspace{4pt}
{\large Construction, bounded results, and failed tests}\par
{\small Solo research discussion note, prepared with AI assistance. Version 3.}

\sect{Abstract}
This note specifies a nonsupersymmetric Spin(10) gauge/scalar construction
with global PQ selection rules and reports bounded calculations. Earlier
locally positive scalar points include a failed Higgs quartic tuning and a
positive-quartic point subsequently rejected by a tree perturbativity screen.
A different exploratory point has positive heavy masses, a positive relaxed
Higgs quartic, a full-field quartic-positivity certificate and a complete
scalar-contact bound below $1/2$. Its formal one-loop singlet pole correction,
however, exceeds its tree squared mass. Gauge unification and complete quantum
consistency remain unresolved. Local anomaly traces and conditional source
identities add necessary checks, not a viability theorem. This is a discussion
draft based on recorded evidence, with no new scientific acceptance decision.

\sect{1. Assumptions and fields}
Work in four spacetime dimensions with gauge group $\Spin(10)$ and global PQ
selection rules. The simply connected group permits the fermionic spinor $16$;
local tensor calculations use the $\mathfrak{so}(10)$ action. A possible quotient
by common gauge/PQ centers has not been fixed, although the ordinary-field
center identification has now been calculated (Section 5). No supersymmetry, explicit soft
PQ violation, or additional scalar representation is assumed.

\begin{center}
\begin{tabular}{lccc}
\toprule Field & Representation/type & $q_{\PQ}$ & Real scalar coordinates\\
\midrule
$\Phi_{ij}$ & real symmetric traceless $54$ & $0$ & $54$\\
$\Sigma_{ijklm}$ & complex self-dual $126$ & $+2$ & $252$\\
$\phi_i$ & complex $10$ & $-2$ & $20$\\
$S$ & complex singlet & $-4$ & $2$\\
$\psi_p$, $p=1,2,3$ & left-handed $16_F$ & $+1$ & ---\\
\bottomrule
\end{tabular}
\end{center}
Indices run from $0$ to $9$, with Euclidean vector metric $\delta_{ij}$ and
$\epsilon_{0\cdots9}=+1$. Fix $*\Sigma=+i\Sigma$; its conjugate is the
$\overline{126}$. The scalar kinetic convention is
\[
\mathcal L_{\rm kin}=
\frac12(D_\mu\Phi_{ij})^*(D^\mu\Phi_{ij})
+\frac{1}{2\,5!}(D_\mu\Sigma_{ijklm})^*(D^\mu\Sigma_{ijklm})
+(D_\mu\phi_i)^*(D^\mu\phi_i)+(\partial_\mu S)^*(\partial^\mu S).
\]
The self-dual factor avoids double counting complementary five-index sectors.
The physical construction uses the usual gauge-curvature kinetic term and
$\mathcal L=\mathcal L_{\rm gauge}+\mathcal L_{\rm kin}+\mathcal L_Y-V$.
The allowed symbolic Yukawa structures are $\psi\psi\phi$ and
$\psi\psi\Sigma^*$, with complex symmetric family matrices and isometrically
normalized intertwiners. Their numerical matrices and right-handed-neutrino
spectrum have not been selected.

The field-content motivation follows Babu and Khan [1]; PQ/axion context is
discussed in [2,3]. The potential below is a separately specified source-derived
variant. Published scalar benchmarks, fermion fits, and unification conclusions
are not imported as results for it. No novelty claim is made.

\newpage
\sect{2. Complete frozen scalar definition}
Write $V=V_2+V_3+V_4$. Repeated tensor indices are summed over all ten values.
Define
\[
\begin{gathered}
N=\frac1{5!}\Sigma_{abcde}\Sigma^*_{abcde},\quad
u=\phi_i^*\phi_i,\quad t=\phi_i\phi_i,\quad
p_n=\Tr(\Phi^n),\ n=2,3,4,\\
K_{ij}=\frac1{4!}\Sigma_{iabcd}\Sigma^*_{jabcd},\qquad
T_{ij}=\frac1{4!}\Sigma_{iabcd}\Sigma_{jabcd},\\
L_{\Phi\Sigma}=\frac1{2\,3!}\Phi_{ij}\Phi_{kl}
\Sigma_{ikabc}\Sigma^*_{jlabc},\qquad
E=\frac1{2!\,2!\,3!}\Sigma_{ijlmn}\Sigma_{ijpqr}\Sigma^*_{lmpqr}\phi_n .
\end{gathered}
\]
$T$ is symmetric traceless. For increasing tuples $I,J$ of length $5-k$,
let $P_k[I,J]=\sum_{K\in\binom{\{0,\ldots,9\}}k}\Sigma_{KI}\Sigma_{KJ}$
and $Q_k=\sum_{I,J}|P_k[I,J]|^2$, for $k=0,1,2$.
Antisymmetric components with repeated indices vanish. If $A,B$ are disjoint,
$\varepsilon(A,B)$ is the parity of their concatenation relative to increasing
order; otherwise it is zero. The fourth real five-form quartic is
\[
\begin{split}
X_{131}={}&\sum_{I,J\in\binom{\{0,\ldots,9\}}4}
\ \sum_{\substack{A\subset I,\ |A|=3\\B\subset J,\ |B|=1}}
\varepsilon(A,I\setminus A)\varepsilon(B,J\setminus B)
\varepsilon(A,B)\varepsilon(I\setminus A,J\setminus B)\\[-2pt]
&\hspace{25pt}\times P_1[I,J]\,
\overline{P_1[\operatorname{sort}(A\cup B),
\operatorname{sort}((I\setminus A)\cup(J\setminus B))]} .
\end{split}
\]
The potential in these normalizations is
\begin{align*}
V_2={}&m_\Phi^2p_2+m_\Sigma^2N+m_\phi^2u+m_S^2|S|^2,\\
V_3={}&\mu_\Phi p_3+\mu_{\Phi\phi}\Phi_{ij}\phi_i^*\phi_j
+[z_6tS^*+\hc],\\
V_4={}&\lambda_{\Phi1}p_2^2+\lambda_{\Phi2}p_4
+\lambda_{\Phi\Sigma1}p_2N+\lambda_{\Phi\Sigma2}L_{\Phi\Sigma}\\
&+\lambda_{\Phi\phi1}p_2u+\lambda_{\Phi\phi2}\phi_i^*\Phi_{ij}\Phi_{jk}\phi_k
+\lambda_{\Phi S}p_2|S|^2\\
&+\lambda_{\Sigma1}Q_0+\lambda_{\Sigma2}Q_1
+\lambda_{\Sigma3}Q_2+\lambda_{\Sigma4}X_{131}\\
&+\lambda_{\Sigma\phi1}Nu+\lambda_{\Sigma\phi2}\phi_i^*K_{ij}\phi_j
+\lambda_{\phi1}u^2+\lambda_{\phi2}|t|^2\\
&+\lambda_{\Sigma S}N|S|^2+\lambda_{\phi S}u|S|^2+\lambda_S|S|^4\\
&+[z_4\Phi_{ij}T_{ij}S+z_K\Phi_{ij}\phi_i\phi_jS^*
+z_\eta E+z_D T_{ij}\phi_i\phi_j+\hc].
\end{align*}
All $m^2,\mu,\lambda$ coefficients are real; $z_6,z_4,z_K,z_\eta,z_D$
are complex. The masses have dimension two, the cubics dimension one, and
the quartics dimension zero. Hermitian conjugation includes the coefficient.
The shorthand $\mu_{\Phi\phi},\lambda_{\phi1,2},\lambda_{\phi S},z_\eta$
corresponds to repository names \pathref{muPhiPhi}, \pathref{lambdaPhiVector1,2},
\pathref{lambdaVectorS}, and \pathref{zEta}.

There are 29 Hermitian coefficient groups: four quadratic, three cubic, and
22 quartic. Their 24 real and five complex coefficients give 34 real directions
before rephasings. Independent $D_5$ character multiplicities give 34 complex
singlet slots before Hermitian pairing; explicit contractions attain that upper
bound. The generic scalar phase quotient leaves three relative coupling phases.
This basis result fixes the action, not a viable vacuum.

\newpage
\sect{3. Scalar results: a failed tuning and a later positive point}
The candidate high-scale background is
\[
\langle\Phi\rangle=\frac{\omega}{\sqrt{60}}
\operatorname{diag}(-2\mathbf1_6,3\mathbf1_4),\quad
\langle\Sigma\rangle=\frac{\sigma}{4\sqrt2}
\bigwedge_{k=0}^{4}(e_{2k+1}+i e_{2k}),\quad
\langle S\rangle=\frac v{\sqrt2},\quad\langle\phi\rangle=0.
\]
The first direction has the intended Pati--Salam/discrete-parity stabilizer;
the combined gauge orbit has 33 broken and 12 unbroken directions, consistent
with the SM Lie algebra. Global topology and the PQ quotient are separate
questions. Tadpoles and Hessians are derived from the same frozen potential.

At the recorded Stage 1 point, $\omega=\sqrt{60}$ in calculation units and
$\sigma/\omega=v/\omega=0.1$. The exact-decimal free coefficients and
tadpole-eliminated masses are specified in \pathref{STAGE1_POINT.json} and its
replay program. All 35 SM-irrep multiplicity blocks were evaluated. The real
Hessian has rank 294 and nullity 34: 33 gauge directions plus one PQ direction.
The remaining eigenvalues are positive in the canonical kinetic metric;
the smallest recorded physical mass squared is approximately $0.08656$.
This is local positivity modulo symmetry, not global boundedness or a
phenomenologically fitted scale. This point has no light Higgs doublet.

The distinct Stage 2 point changes
$\operatorname{Re}z_6,\operatorname{Re}z_K,\operatorname{Re}z_\eta$ to $0.3$
in the stored dimensionless convention and tunes
$m_\phi^2/\omega^2\simeq-0.1337864654$ to the exact algebraic root recorded in
\pathref{STAGE2_POINT.json}. It has rank 290 and nullity 38, including four
real coordinates of one massless complex SM Higgs doublet. All physical heavy
quadratic modes remain positive. This statement alone is insufficient for
local stability along the tuned flat direction.

\sect{Decisive additional test}
For a canonically normalized neutral field $h$, $H^\dagger H=h^2/2$, expand
around the tuned point with physical heavy coordinates $X_A$:
\[
V=V_0+a_{\rm direct}h^4+h^2j_AX_A
+\tfrac12X_A\mathcal H_{AB}X_B+\cdots .
\]
Minimization on the positive heavy subspace gives
$X=-\mathcal H^+j h^2+\cdots$, and hence
\[
\lambda_{\rm eff}=4\left(a_{\rm direct}
-\tfrac12j^T\mathcal H^+j\right),\qquad
V_{\rm eff}=\lambda_{\rm eff}(H^\dagger H)^2+\cdots .
\]
The recorded parent-action evaluation gives
\begin{center}
\begin{tabular}{lr}
\toprule Quantity in $(H^\dagger H)^2$ normalization & Value\\
\midrule Fixed-heavy direct quartic & $+0.19910$\\
Heavy relaxation subtraction & $56.68893$\\
Relaxed effective quartic & $\boldsymbol{-56.48984}$\\
\bottomrule
\end{tabular}
\end{center}
The negative sign makes the Stage 2 point a failed local scalar minimum at
tree level. It does not invalidate the separate Stage 1 result or prove that
every parameter point fails. A direct parent-potential expansion along the
heavy-field valley reproduces the quartic; heavy-coordinate shear and complex
Higgs-phase controls preserve it. These extraction routes share the frozen
tensor implementation. No global-minimum or loop-corrected stabilization claim
is made.

\sect{A distinct positive-quartic point}
The later recorded point in \pathref{canonical_so10_positive_higgs/POINT.json}
retains the Stage 1 VEVs but uses $z_6/\omega=z_K=0.03$, $z_\eta=0.3$,
$\lambda_{\Phi\phi1}=-0.002$ and an exact algebraic tuning
$m_\phi^2/\omega^2\simeq0.007302047861$. Its real Hessian has rank 290 and
nullity 38, with positive heavy modes. The direct quartic is $0.19845235$;
heavy relaxation subtracts $0.09359411$, leaving
\[
\boxed{\lambda_{\rm eff}=+0.10485824.}
\]
The point's SHA-256 is \texttt{476645438ea8a7ada42fea90e80174b772}\newline
\texttt{d45ce3e729d013d134e9848950c816}.
Direct-parent and valley-polynomial routes agree with shared tensor inputs.
This remains a local tree-level result. Subsequent calculation at this SAME
point gives a physical coherent scalar channel with $|a_0|=0.76809032>1/2$
in the fixed-input high-energy tree limit [7]. Its vector-exchange currents
vanish in that channel. The point therefore fails this perturbativity screen;
positive local masses and quartic did not establish a controlled benchmark.
This failure does not exclude the construction. The new exploratory point
in Section 6 has different inputs and must be assessed separately.

\newpage
\sect{4. Gauge matching: a measured mismatch, an incomplete verdict}
The recorded direct one-loop calculation matches the full parent to the SM
in a common $\overline{\rm MS}$ background-field Feynman-gauge convention,
with GUT-normalized $g_1$. It includes 290 heavy real scalar directions and
33 massive vector/Goldstone/ghost sets from the later positive-quartic point,
not the earlier failed tuning. The threshold ledger binds that point's hash.
This point subsequently fails the scalar perturbativity screen. These
residuals have not been recomputed for the new exploratory point in Section 6.

Define $\alpha_i=g_i^2/(4\pi)$ and one-loop threshold functions $\lambda_i$ by
\[
\alpha_i^{-1}(\mu)=\alpha_{10}^{-1}(\mu)-\frac{\lambda_i(\mu)}{12\pi}.
\]
The complete state accounting gives
\[
\begin{gathered}
b_{\rm SM}=\left(\frac{41}{10},-\frac{19}{6},-7\right),\quad
b_S=\left(\frac{377}{30},\frac{77}{6},\frac{79}{6}\right),\quad
T_V=(8,6,5),\\
b_{{\rm SM},i}+b_{S,i}-\frac72 T_{V,i}=-\frac{34}{3}=b_{10},\qquad
\frac{d\lambda_i}{d\ln\mu}=6(b_{{\rm SM},i}-b_{10}).
\end{gathered}
\]
These are exact beta-jump and one-loop matching-scale checks. At
$\mu=\omega$, the frozen inverse-coupling inputs are
$(59.0100,29.5800,8.47458)$ at $M_Z$. The one-loop equations are
\[
\alpha_i^{-1}(M_Z)=\frac{4\pi}{g_{10}^2}
-\frac{\lambda_i^S+\lambda_i^V(g_{10})}{12\pi}
+\frac{b_i^{\rm SM}}{2\pi}\ln\frac{\omega}{M_Z}.
\]
They have no common perturbative solution for the frozen inputs. The recorded
least-squares residuals, in inverse-coupling units and channel order $(1,2,3)$,
are:
\begin{center}
\begin{tabular}{lrrrr}
\toprule Completed approximation & $r_1$ & $r_2$ & $r_3$ & $\max|r_i|$\\
\midrule One-loop running and thresholds & $0.36615$ & $-1.05980$ & $0.69389$ & $1.05980$\\
Two-loop SM running, one-loop thresholds & $0.26091$ & $-0.74548$ & $0.48479$ & $0.74548$\\
\bottomrule
\end{tabular}
\end{center}
The latter uses the SM running framework [5] and frozen reference boundaries
from [6], including top, bottom, and tau insertions with their one-loop
evolution. A separate inverse-coupling evolution with a different adaptive
integrator reproduces the residual to better than $2\times10^{-8}$ under
the frozen equations. Matching-scale variation $\mu/\omega\in[0.5,2]$
leaves maximum residuals from $0.74172$ to $0.74931$.

\sect{What the missing calculation must supply}
The exact two-loop RG logarithmic template has been derived, but its finite
model-dependent heavy constants have not. The published general matching
framework [4] provides methodology and a restricted comparator, not the finite
answer for these mixed, nondegenerate fields and multiple VEVs. The complete
Yukawa-dependent threshold also requires numerical family matrices and
right-handed-neutrino masses.

Small scale variation is not a rigorous bound on unknown finite constants.
Setting them to zero or freely tuning them would not compute a prediction.
Consequently the formal gauge verdict remains unresolved: neither unification
success nor exclusion has been established. The diagnostic fitted
$g_{10}\simeq0.5231$ and $\omega\simeq1.64\times10^{15}\,\mathrm{GeV}$
are not physical benchmark parameters. No replacement spectrum or flavor fit
is introduced in this note.

\newpage
\sect{5. Symmetry and quantum evidence: scope matters}
A separately named covariant physical gauge/scalar candidate satisfies exact
classical local gauge identities through factorized representation checks.
The recorded primary and tensor-functor replay compare the 328-real scalar
representation action and all 45 generator directions. The usual
$\delta(Dx)=-g_{10}\rho(\omega)Dx$ identity, invariant kinetic metric,
curvature identity, and invariant potential establish the classical statement.
This qualification is inactive; it does not replace a promoted component AST
or complete a quantum calculation.

Earlier literal component-action failures are retained. One preserved
two-field witness has nonzero gauge variation
\[
\delta\mathcal L_{\rm old}=
\frac{i g_{10}}{\sqrt2}(\partial_\mu C^4)
(Y\partial^\mu X-X\partial^\mu Y),
\]
with nonzero ghost Euler derivative. The later covariant candidate fixes the
kinetic structure by derivation. A failed encoded candidate is evidence about
that representation of the action, not a theorem against the gauge group.

Separately, the recorded one-loop parent scalar-pole projection has exact
rank $4+4+26=34$ and zero out-of-basis residual in its tested operator space.
Hermiticity and PQ selection-rule audits also pass within that scope. This
bounded radiative-closure result is not a complete anomaly or quantum
consistency theorem; it does not evaluate finite two-loop thresholds.

In the gauge-fixing/source investigation, a preserved normalization candidate
misses a selected auxiliary--antighost--ghost cubic pole. At one recorded
real-component witness, in common
$\hbar g_{10}\xi^2/(32\pi^2\epsilon)$ units, the coefficients are
\[
G_{\rm actual}=\frac{48\sqrt{10}}{25},\quad
G_{\rm candidate}=\frac{8\sqrt{10}}5,\quad
G_{\rm actual}-G_{\rm candidate}=\frac{8\sqrt{10}}{25}\ne0.
\]
Later calculations retain this failed candidate while deriving actual
additional source, auxiliary and ghost coefficients and checking selected
local Ward/Slavnov relations. These are scoped operator results. They do not
establish the full source-free Slavnov identity or a quantum-equivalence theorem.

The latest selected background test verifies that a specified broken-antighost
operation commutes with the actual local residual-gauge transformation on all
1334 classical field/source/background columns. In its declared polynomial
ring, a 642-column Ward/Euler calculation proves uniqueness of a background
extension at coupling orders zero and one, conditional on locality and inherited
covariance of the actual logarithmic UV pole. A nonlocal transverse
counterexample shows why locality is essential. Subsequent tests verify weighted principal invertibility on all 643 fields
and construct one explicit local covariant reference operator, recovering all
1597 flat entries. This supplies a local-symbol member, without identifying a
unique historical completion or adopting a physical scheme. The actual product
of heat-damped propagators still lacks a proved UV/subtraction correspondence;
the full background logarithmic identity remains conditional.
Regulated insertions, products/contacts, measure and infrared limits remain open.
The formal ledger is still at 34 of 39 items with the quantum gate blocked.

\sect{Gauge anomaly and PQ structure}
For the three left-handed $16_F$ multiplets, the local four-dimensional
perturbative gauge-anomaly tensor and linear gauge traces vanish. A weight
polynomial route with conjugacy/polarization and a Clifford-matrix route
evaluate the established representation criterion [9], including all 16,215
symmetric cubic slots. This does not construct a regulated matter determinant
or settle global anomalies, bundles or the full quantum measure.

For the declared PQ charges the convention $N=2\sum_f q_f T_c(f)$ and
$E=2\sum_f q_f Q_f^2$ gives
\[
N=12,\qquad E=32,\qquad E/N=8/3.
\]
The common Spin(10) center identifies an ordinary-field PQ quarter-turn;
nonzero gauge-singlet $S$ fixes the vacuum-orbit period to $\alpha=\pi/2$.
The conventional QCD potential would consequently give $N_{\rm DW}=3$,
\emph{conditional} on realistic nonzero quark masses/QCD susceptibility, this
ordinary orbit and no additional identification or explicit PQ breaking [3].
The electroweak-unbroken reference does not supply that QCD vacuum. This is
neither a physical axion mass/abundance prediction nor a full global completion.

\newpage
\sect{6. A new scalar witness and its quantum-control warning}
The different rational fixture in
\pathref{canonical_so10_coherent_joint_witness/POINT.json}
uses the same field content, action and VEV shape. Its frozen reduced inputs
are $p_1=1/100$, $p_2=1/50$, $x=3/25$, $s=1/1000$, $z=1/500$,
$D=1/50$, $t=23/1250$, $\Delta=0$, $\eta=1/10000$; the dependent
coefficients and larger simple weak root are fixed by the recorded action
maps. It is an exploratory fixture, not an adopted physical benchmark.
Do not combine its results with the old point's spectrum, thresholds or loops.

Exact specialization of the inherited 12 spectral types (35 blocks, 328 real
fields) gives 290 positive directions and 38 symmetry/tuning zeros. The
actual weak kernel and all positive neutral relaxation modes give
$\lambda_{\rm eff}\simeq0.1074374161>0$ at tree level. This uses a new
point-specific light direction, not a transported earlier quartic.

\sect{Complete classical coverage at this point}
The complete symmetric scalar pair space has 53,956 states. In the contact
convention of [7], the recorded structural Gram and block-norm bounds give
\[
\|a_0^{\rm contact}\|\le
\frac{35439836953}{24000000000\pi}
\simeq0.47003544<\frac12.
\]
This is a conservative upper bound, not a largest eigenvalue or complete
gauged scattering amplitude. Massless-gauge exchange requires its own
infrared-safe treatment; no gauge coupling is silently set to zero.

With $X=\Tr\Phi^2$, $n=N_\Sigma$, $u=\phi^\dagger\phi$, $v=|S|^2$,
global invariant inequalities and explicit completion of the complex mixed
terms give a rational positive-definite four-norm lower matrix. Its exact
minors and a separate rational solver yield, for all 328 real fields,
\[
V_4\ge c(X+n+2u+2v)^2,\qquad
c=\frac{824404010711173266247333949}
{1720579893948342670590110120750}
\simeq4.79143115\times10^{-4}>0.
\]
This is global quartic positivity, rather than sampled positive rays. It
ensures growth at large fields for this potential. It does not prove that the
chosen stationary vacuum is deepest or has a sufficient tunneling lifetime.
The old failed points and closed parameter families remain failed in their
qualified scopes; these certificates do not transfer to other inputs.

\sect{A complete selected one-loop pole coefficient is large}
At the same fixture, $Q^2=\omega^2=60$ and the declared
$\overline{\rm MS}$ stationary-branch expansion, the isolated singlet has
tree squared mass $6/5$. Scalar seagulls and momentum-dependent bubbles,
including the physical PQ/Higgs zero modes, give a one-particle-irreducible
contribution $-1.36799138-0.23557167i$. The required vacuum-source insertion
adds $+3.03144395$, so [8]
\[
\delta m^2_{S,\rm pole}=1.66345257-0.23557167i,\qquad
\frac{|\delta m^2_{S,\rm pole}|}{m^2_{S,\rm tree}}\simeq1.40004.
\]
Field-content and projector checks justify this particular one-loop component.
Direct singlet gauge/Yukawa vertices vanish here. The tree singlet plane is
isolated and the relevant other tadpole multipliers vanish. This does not
make every quantum-vacuum component complete.
An exact local RG/log check retains generated $\Phi$--$S$ and $\Sigma$--$S$
portals [10] and cancels the scale derivative, but does not reduce the large
correction. A small first-order correction is unsupported in this expansion.
No reliable physical pole mass/width, resummed solution or universal scheme
failure is inferred. Classical feasibility and controlled quantum viability
remain distinct.

\newpage
\sect{7. Evidence and reproduction map}
This note summarizes recorded evidence at source commit
\begin{center}\texttt{22709b33f20484efe16eb4c3592d690c8e0650f3}.\end{center}
The source scientific authority remains \pathref{project_state.json}, version
113, unchanged by publication. The original Version 2 note and two base source
archives remain preserved. The Version 2 \emph{source update} overlays those
archives; the reproduction guide specifies extraction, deleted paths and
separate historical-policy/external-control namespaces. Manifests authenticate
bytes without performing scientific acceptance. Original Git history is not
included; historical custody exit 2 with seven unresolved references is not PASS.

Evidence paths below are relative to the assembled source tree. Each new row
includes its prospective \pathref{RECORD.md}, result, actual validation receipt,
frozen inputs, sources and complete successful/failed captures.

{\small
\begin{tabular}{p{0.21\linewidth}p{0.73\linewidth}}
\toprule Claim & Evidence locator under \pathref{calculations/}\\
\midrule
Scalar action & \pathref{canonical_so10_scalar_reconstruction/RESULT_V2.md}\\[3pt]
Original masses & \pathref{canonical_so10_scalar_benchmark/RESULT.md}\\[3pt]
Failed / positive quartics & \pathref{canonical_so10_light_higgs_quartic/RESULT.md} and\newline
\pathref{canonical_so10_positive_higgs/RESULT.md}\\[3pt]
Old point contact failure & \pathref{canonical_so10_radial_input_control/RESULT.md}\\[3pt]
New point and local tests & \pathref{canonical_so10_coherent_joint_witness/RESULT.md}\\[3pt]
Whole contact bound & \pathref{canonical_so10_complete_contact_bound/RESULT.md}\\[3pt]
Global quartic bound & \pathref{canonical_so10_finite_full_bfb/RESULT.md}\\[3pt]
Singlet pole warning & \pathref{canonical_so10_singlet_pole_control/RESULT.md}\\[3pt]
Pole RG/log check & \pathref{canonical_so10_singlet_pole_rg_check/RESULT.md}\\[3pt]
Gauge residuals & \pathref{canonical_so10_direct_gauge_matching/RESULT.md} and\newline
\pathref{canonical_so10_direct_gauge_next_order/RESULT.md}\\[3pt]
Missing finite boundary & \pathref{canonical_so10_direct_two_loop_heavy_threshold/RESULT.md}\\[3pt]
Gauge / PQ traces & \pathref{canonical_so10_local_gauge_anomaly/RESULT.md} and\newline
\pathref{canonical_so10_pq_anomaly_screen/RESULT.md}\\[3pt]
Conditional background & \pathref{canonical_so10_ub12_background_locality/RESULT.md}\\
\bottomrule
\end{tabular}}

The scalar-pole closure record is under
\pathref{tools/two_loop_qft_compiler_v1/LAYER5B_M05_M08_RESULT.md}.
The preserved classical-action and auxiliary-normalization records are in
its \pathref{m08_independent_replay_package_v1/} subdirectory:
\pathref{A10_PHYSICAL_ACTION_RECONCILIATION_RESULT.md} and
\pathref{A10_E5_AUXILIARY_VERTEX_RESULT.md}.
The larger source archive retains intermediate exclusions, corrections and
bounded source/ghost calculations; archive inclusion does not promote every
historical claim. The completed weighted-reference and local-witness records
are also included; work after the frozen checkpoint is excluded.

\sect{What publication validation does}
The package verifier checks every released file and archive member. The
publication receipt separately authenticates the selected scientific receipts
and their immutable inputs, distinguishes historical policy versions and
records executed short checks. No expensive Hessian, contact, singlet-pole or
full quantum calculation is repeated for this release. Original paths in
external control scripts may require explicit relocation in a disposable
environment; flat export is not an automatic historical audit replay.

\newpage
\sect{8. Remaining obligations and AI assistance}
No physically viable GUT or experimental agreement is established. The new
point has a full classical quartic bound, but the preferred vacuum, tunneling,
loop control, hierarchy/naturalness and complete gauged scattering remain
open. Full quantum restoration/measure/contact/IR obligations, finite heavy
thresholds, realistic matter/neutrino inputs, complete axion cosmology and
proton-decay phenomenology also remain open. The old gauge residuals are
diagnostics at their own frozen inputs and do not predict the new point.

Useful discussion questions are whether the new point's large singlet
correction admits a justified, controlled physical-input expansion, and what
model-specific locality and finite-threshold calculations are needed before
extending the quantum and gauge conclusions. Missing contributions are not
assumed to rescue any failed result.

LLMs assisted derivations, code, verification and writing. Scientific direction
and acceptance belong to the human author. The records demonstrate bounded
symbolic/numerical calculation, structural reductions and cross-checking;
they also preserve failed interpolation, assertion, normalization and
verification attempts. Corrected agreement cannot validate a shared physical
assumption. Shared action/representation inputs, inherited proof caches,
some tensor/arithmetic code, Python/SymPy and the same analyst limit
independence. Selected work uses NumPy, SciPy, mpmath, CVXPY/Clarabel and
Julia. Multiple AI routes are not independent scientists. No external expert
verification, whole-model proof-assistant certification or general reliability
claim is made.

\sect{References}
{\small\setlength{\parskip}{3pt}
[1] K. S. Babu and S. Khan, \emph{A Minimal Non-Supersymmetric SO(10) Model:
Gauge Coupling Unification, Proton Decay and Fermion masses},
Phys. Rev. D \textbf{92}, 075018 (2015).
\href{https://arxiv.org/abs/1507.06712}{arXiv:1507.06712}.

[2] R. D. Peccei and H. R. Quinn, \emph{CP Conservation in the Presence of
Pseudoparticles}, Phys. Rev. Lett. \textbf{38}, 1440 (1977).
\href{https://doi.org/10.1103/PhysRevLett.38.1440}{doi:10.1103/PhysRevLett.38.1440}.

[3] A. Ernst, A. Ringwald and C. Tamarit, \emph{Axion Predictions in
SO(10) $\times$ U(1)$_{\PQ}$ Models}, JHEP \textbf{02} (2018) 103.
\newline\href{https://arxiv.org/abs/1801.04906}{arXiv:1801.04906}.

[4] W. Martens, \emph{Towards a Two-Loop Matching of Gauge Couplings in
Grand Unified Theories}, JHEP \textbf{01} (2011) 104.
\href{https://arxiv.org/abs/1011.2927}{arXiv:1011.2927}.

[5] M. Luo and Y. Xiao, \emph{Two-loop Renormalization Group Equations in
the Standard Model}, Phys. Rev. Lett. \textbf{90}, 011601 (2003).
\href{https://arxiv.org/abs/hep-ph/0207271}{arXiv:hep-ph/0207271}.

[6] D. Buttazzo et al., \emph{Investigating the near-criticality of the Higgs
boson}, JHEP \textbf{12} (2013) 089.
\href{https://arxiv.org/abs/1307.3536}{arXiv:1307.3536}.

[7] M. D. Goodsell and F. Staub, \emph{Unitarity constraints on general scalar
couplings with SARAH}, Eur. Phys. J. C \textbf{78}, 649 (2018).
\href{https://arxiv.org/abs/1805.07306}{arXiv:1805.07306}.

[8] S. P. Martin, \emph{Two-loop scalar self-energies in a general
renormalizable theory at leading order in gauge couplings}, Phys. Rev. D
\textbf{70}, 016005 (2004).
\href{https://arxiv.org/abs/hep-ph/0312092}{arXiv:hep-ph/0312092}.

[9] A. Bilal, \emph{Lectures on Anomalies} (2008).
\href{https://arxiv.org/abs/0802.0634}{arXiv:0802.0634}.

[10] M. Luo, H. Wang and Y. Xiao, \emph{Two-loop Renormalization Group
Equations in General Gauge Field Theories}, Phys. Rev. D \textbf{67},
065019 (2003). \href{https://arxiv.org/abs/hep-ph/0211440}{arXiv:hep-ph/0211440}.

}
\end{document}
